\documentclass[letterpaper,12pt]{article}
%\usepackage{harvard}
%\usepackage{subfigure}
%\usepackage[active]{srcltx}
\usepackage{graphicx}
\usepackage{pdflscape}
\usepackage{amsmath}
\usepackage{amsthm}
\usepackage{lscape}
\usepackage{enumerate}
\usepackage{float}
\usepackage{bm}
\usepackage{grffile}
\usepackage{relsize}
\usepackage{tablefootnote}
\usepackage{footnote}
\usepackage{sectsty}
\usepackage{xcolor}
\usepackage{caption}
\usepackage{color}
\usepackage{natbib}
\usepackage{multirow}
\usepackage{bbm}
\usepackage[T1]{fontenc}
\usepackage[multiple]{footmisc}
\usepackage{appendix}
\usepackage{times}
\usepackage{relsize}
\usepackage{caption}
\usepackage{newtxtext,newtxmath}
\usepackage{enumitem}


%%%%%%%%%%%%%%%%%%%%%%%%%%%%%%%%%%%%%%%%%%%%%%%%%%%%%%%%%%%%
%%%   Title Page Information
%%%%%%%%%%%%%%%%%%%%%%%%%%%%%%%%%%%%%%%%%%%%%%%%%%%%%%%%%%%%
\title{\vspace{-.5in} ECON712: Problem Set 1}
\author{Leandro Freylejer}
\date{Due date: February 13th, 2026 before 4pm }

%%%%%%%%%%%%%%%%%%%%%%%%%%%%%%%%%%%%%%%%%%%%%%%%%%%%%%%%%%%%
%%%   Double and Single Space Commands
%%%%%%%%%%%%%%%%%%%%%%%%%%%%%%%%%%%%%%%%%%%%%%%%%%%%%%%%%%%%
\newlength{\defbaselineskip}
\setlength{\defbaselineskip}{\baselineskip}
\newcommand{\setlinespacing}[1]{\setlength{\baselineskip}{#1 \defbaselineskip}}
\newcommand{\doublespacing}{\setlength{\baselineskip}{1.3 \defbaselineskip}}
\newcommand{\singlespacing}{\setlength{\baselineskip}{1\defbaselineskip}}
\renewcommand{\thesubsection}{\arabic{subsection}}
\newtheoremstyle{dotlessS}{}{}{\itshape}{}{\color{dg}\bfseries}{}{ }{}
\theoremstyle{dotlessS}
\newtheorem{theorem}{Theorem}
\newtheorem{lemma}{Lemma}
\newtheorem{proposition}{Proposition}
\newtheorem{assumption}[theorem]{Assumption}
\newtheorem{corollary}{Corollary}
\newtheorem{definition}{Definition}
\newenvironment{example}[1][Example]{\begin{trivlist}
\item[\hskip \labelsep {\bfseries #1}]}{\end{trivlist}}
\newenvironment{remark}[1][Remark]{\begin{trivlist}
\item[\hskip \labelsep {\bfseries #1}]}{\end{trivlist}}

\newcommand\Pitem{%
  \addtocounter{enumi}{-1}%
  \renewcommand\theenumi{\arabic{\enumi}'}%
  \item%
  \renewcommand\theenumi{\arabic{enumi}}%
}

%%%%%%%%%%%%%%%%%%%%%%%%%%%%%%%%%%%%%%%%%%%%%%%%%%%%%%%%%%%%
%%%   Set Margins Here
%%%%%%%%%%%%%%%%%%%%%%%%%%%%%%%%%%%%%%%%%%%%%%%%%%%%%%%%%
\textwidth=6.5in
\textheight=9in
\oddsidemargin=0.10in
\evensidemargin=0.10in
\topmargin=-0.5in
%\parskip=3pt

%%%%%%%%%%%%%%%%%%%%%%%%%%%%%%%%%%%%%%%%%%%%%%%%%%%%%%%%%%%%%
%%%   Actual Document Begins Here
%%%%%%%%%%%%%%%%%%%%%%%%%%%%%%%%%%%%%%%%%%%%%%%%%%%%%%%%%%%%%
\begin{document}
\pagenumbering{arabic}

\maketitle
\vspace{-.3in}
\textbf{Instructions:} There are $5$ questions for a total of $52$ points. Answer all questions. To receive full marks, write your answers as complete and clear as possible. For the empirical question, please attach your code so I can replicate your results (if using STATA please submit log files). 

\section*{Question I: Conditional Expectation Function}
Suppose that the random variables $Y$ and $X$ only take the value $0$ and $1$ and they have the following joint probability distribution: 
\begin{table}[h]
\centering
\begin{tabular}{|c| c |c|}
\hline
         & X=0 & X=1 \\
         \hline
Y=0  & 0.1 & 0.2 \\
Y=1  & 0.4 & 0.3 \\
\hline
\end{tabular}
\label{tab:three_by_three}
\end{table}
\\
Using this information answer the following questions [Hint: to answer the questions below you will need to use the fact that both the distribution of $Y$ and  conditional distribution of $Y|X$ are Bernoulli]:
\begin{enumerate}[label=(\alph*)]
\item (5 points) Find $E[Y|X=x]$, $Var[Y|X=x]$ for $x=0$ and $x=1$. Provide an interpretation of  $E[Y|X=x]$ in terms of probabilities.  
\item (5 points) Create a conditional expectation function $m(X)$, such that $Y=m(X)+u$ and $E[u|X=x]=0$ for $x=0,1$. Solve for $Var[u]$.
\item (5 points) Suppose that a researcher uses the method-of-moment estimator (MM), $E[XU]=0$ on a random sample of $n$ individuals to estimate the parameters of the following regression:
\begin{eqnarray*}
y_i=\hat{\beta_0}+\hat{\beta_1} x_i+\hat{u}_i
\end{eqnarray*}
Interpret  $\hat{\beta_0}$ and $\hat{\beta_1}$ as functions of sample moments. Given your previous results, discuss any violation of the Gauss-Markov (Classical) assumptions that results in MM/OLS being the best linear unbiased estimator [Hint: the issue is related to $Var[Y|X=x]$ that you calculated in part (a)] .       
\item (5 points) Demonstrate the Law of Iterated Expectations explicitly by showing $E[E[X|Y]]=E[X]$
\end{enumerate}
\newpage


\section*{Question II: Linear Transformation of Dependent Variable (7 points)}
Consider the linear regression model:
\begin{eqnarray*}
\bm{y}=\bm{X} \: \bm{\hat{\beta}}+\bm{\hat{u}}
\end{eqnarray*}
where $\bm{X}$ is an $nxk$ matrix of regressors that includes a column of ones, and $\bm{\hat{\beta}}$ is the MM/OLS estimator based on $n$ observations.\\
Define a transformed dependent variable $\bm{z}$ by
\begin{eqnarray*}
\begin{array}{lr}
(i) \: \:  \bm{z}=\bm{y}-c \: \bm{\iota} & (ii) \: \:  \bm{z}=\bm{A} \: \bm{y}
\end{array}
\end{eqnarray*}
where $\bm{\iota}$ is a $nx1$ vector of ones, $c$ is a scalar constant, and $\bm{A}$ is an $nxn$ diagonal matrix with all diagonal elements equal to $c$. Let $\bm{\hat{\alpha}}$ denote the MM/OLS estimator obtained by regression $\bm{z}$ on $\bm{X}$. Derive the relationship between $\bm{\hat{\beta}}$ and $\bm{\hat{\alpha}}$ in cases $(i)$ and $(ii)$.
    

\section*{Question III: Linear Transformation of Regressors (7 points)}
Consider two regressions:
\begin{eqnarray*}
\begin{array}{ll}
\bm{y}=\beta_1 \: \bm{x}_1+\beta_2 \: \bm{x}_2+\beta_3 \: \bm{x}_3+\bm{u} &\text{, and}  \\[1.2ex]
\bm{y}=\beta_1 \: \bm{z}_1+\beta_2 \: \bm{z}_2+\beta_3 \: \bm{z}_3+\bm{u} &
\end{array}
\end{eqnarray*}
where $\bm{z}_1=\bm{x}_1-2 \: \bm{x}_2$, $\bm{z}_2=\bm{x}_2+4 \: \bm{x}_3$, and $\bm{z}_3=2 \: \bm{x}_1-3 \: \bm{x}_2+5 \: \bm{x}_3$. Let $\bm{X}=[\bm{x}_1 \: \: \bm{x}_2 \: \: \bm{x}_3]$ and $\bm{Z}=[\bm{z}_1 \: \: \bm{z}_2 \: \: \bm{z}_3]$. Show that the columns of $\bm{Z}$ can be expressed as linear combinations of the columns of $\bm{X}$, that is $\bm{Z}=\bm{X \: A}$, for some $3x3$ matrix $\bm{A}$. Find the elements of matrix $\bm{A}$.  \\
\\
Show that matrix $\bm{A}$ is invertible by showing that the columns of $\bm{X}$ are linear combinations of the columns of of $\bm{Z}$. Give the elements of $\bm{A}^{-1}$. Show that the two regressions give the same fitted values and residuals. \\
\\
Precisely how is the OLS estimates of $\bm{\hat{\beta}}_1$ related to $\bm{\hat{\alpha}}_i$ for $i=1,2,3$ ?  Precisely how is the OLS estimates of $\bm{\hat{\alpha}}_1$ related to $\bm{\hat{\beta}}_i$ for $i=1,2,3$ 

\section*{Question IV: Empirical Application}
For this question you will need to download data from the Labour Force Survey (LFS). The microdata is provided to Canadian Universities through the DLI (Data Liberation Initiative). To download the data go to the Weller Library website and click on Databases. Search for \textit{Abacus} and access the database. In the Database you need to select the University of Northern British Columbia and create an account with your institutional email (if you are on campus the account is created automatically with your first login). \\
 \\
Within Abacus Data Network search for Labour Force Survey 2025. The LFS is a monthly dataset, so you will need to download each month separately and then combine them, using the \textit{append} command in STATA, to obtain all the data for the year 2025. To download data choose the format `.tab' so you can upload the data set into STATA (there is a lot of documentation online that will guide you through the process of importing tab delimited data into STATA). You should also download the data dictionary to learn about the definition of each variable.\\
\\
In this question you will be tasked with studying the returns to completing University defined as attaining a four-year bachelor's degree or higher. Using this dataset, answer the following questions:
\begin{enumerate}[label=(\alph*)]  
\item (5 points) The first task is to create a table in the style of Table 2.2.1 in \textit{Mostly Harmless Econometrics}. The first two columns of the table (after the variable names) should contain averages conditional on attaining and not attaining a bachelor's degree for $ln(hourly_wages)$, $tenure$ (with current employment), $age$, $female$, $immigrant$. For the last two variables, since they are dummy variables, you will calculate the share of females and immigrants among workers with and without a bachelor's degree. The last column should contain the differences between these variables across the two groups of workers (with and without a bachelor's degree). In parenthesis below the calculated differences, you will need to show the standard error of the difference. Comment on the differences across these two groups of workers. [Hint: to create the last column, think about which regression you will use to obtain the standard error of the difference.]    
\item (5 points) The second task is create a table in the style of Table 2.2.2  in \textit{Mostly Harmless Econometrics}. In particular, you will need to show the estimates of the following regression:
\begin{eqnarray*}
y_i=\beta_0+\beta_1 \: Univ_i+\bm{x}'_i \bm{\alpha}+u_i
\end{eqnarray*}
where $y_i$ is the log of hourly wage, $Univ_i$ is a dummy equal to one if individual $i$ completed at least a bachelor's degree, and $\bm{x}_i$ is a column vector that contains $age$, $age^2$, a dummy for $female$, $tenure$ and a dummy for recent $immigrant$. To replicate the table start with a regression of $y_i$ on a constant and $Univ_i$ and add the other variables on-by-one (age and age$\_$squared should be added together). Make sure you display standard errors in parenthesis below the estimates. Comment on the differences in the estimate of the return to university. Is this estimate causal? Explain.
\item (3 points) Suppose that $u_i=\epsilon_i+\rho \: a_i$ where $a_i$ denotes person $i$'s ability. Also, assume that $E[\epsilon|Univ, \bm{x}]=0$ and that $Cov(a,Univ)>0$. Assuming homogenous effects of attending university on income, express the selection bias of the regression in part (b) in estimating $\beta_1$ as a function of parameters and conditional expectation of $a$.         
\end{enumerate}

\section*{Question V: Projections (5 points)}
Show that $\bm{M_X \: M_{X_1}}=\bm{M_{X_1}\: M_X}=\bm{M_X}$, where $\bm{X}=[\bm{X_1} \: \bm{X_2}]$. Interpret this result

\end{document}